\chapter{Συμπεράσματα και Μελλοντικές Επεκτάσεις}
\label{ch:conclusions}

Η παρούσα διπλωματική εργασία πραγματεύτηκε τη σχεδίαση και υλοποίηση ενός
πειραματικού \textlatin{Data Space} με χρήση τεχνολογιών
\textlatin{FIWARE} και του πλαισίου \textlatin{i4Trust}/\textlatin{iSHARE}.
Στόχος ήταν η δημιουργία ενός λειτουργικού πρωτοτύπου που να αποδεικνύει τον
κύκλο εκχώρησης πρόσβασης σε δεδομένα μεταξύ δύο ανεξάρτητων οργανισμών.

% ------------------------------------------------------------------
\section{Κύρια Συμπεράσματα}
\label{sec:results}
% ------------------------------------------------------------------

Από την υλοποίηση και την αξιολόγηση του συστήματος προέκυψαν τα ακόλουθα
συμπεράσματα:

Η υλοποίηση αποδεικνύει ότι η δημιουργία ενός λειτουργικού \textlatin{Data
Space} είναι εφικτή αποκλειστικά με ανοιχτά εργαλεία και πρότυπα. Τα
συστατικά που χρησιμοποιήθηκαν, \textlatin{Kong}, \textlatin{Keyrock},
\textlatin{Orion-LD} και \textlatin{iSHARE Satellite}, είναι ελεύθερα
διαθέσιμα και καλά τεκμηριωμένα, καθιστώντας την προσέγγιση αυτή
προσιτή για οργανισμούς κάθε μεγέθους.

Ο μηχανισμός εκχώρησης εξουσιοδότησης του \textlatin{iSHARE} αποδείχθηκε
αποτελεσματικός στην υλοποίηση της ψηφιακής κυριαρχίας. Ο πάροχος δεδομένων
διατηρεί πλήρη έλεγχο: μπορεί να εκδίδει, να τροποποιεί και να ανακαλεί
πολιτικές πρόσβασης ανά πάσα στιγμή, χωρίς να απαιτείται η επέμβαση τρίτου
φορέα. Ο καταναλωτής αποκτά πρόσβαση αποκλειστικά στα δεδομένα για τα
οποία έχει εκδοθεί ρητή πολιτική και μόνο για αυτά.

Ο διαχωρισμός μεταξύ Σημείου Επιβολής Πολιτικής (\textlatin{Kong}) και
Σημείου Λήψης Απόφασης (\textlatin{DSBA-PDP}) αποδείχθηκε ιδιαίτερα
χρήσιμος από αρχιτεκτονική σκοπιά. Καθένα από τα δύο συστατικά επιτελεί
ξεκάθαρα οριοθετημένο ρόλο, καθιστώντας το σύστημα επεκτάσιμο και
ευκολότερα συντηρήσιμο.

Η υλοποίηση ανέδειξε αρκετές πρακτικές προκλήσεις που δεν είναι άμεσα
ορατές από τη θεωρητική μελέτη των προδιαγραφών:

\begin{itemize}
    \item Συμβατότητα κρυπτογραφικών μορφών: Το \textlatin{Kong}
    απαιτεί ιδιωτικό κλειδί σε μορφή \textlatin{PKCS\#1}, ενώ τα
    περισσότερα εργαλεία παράγουν κλειδιά σε μορφή \textlatin{PKCS\#8}.
    \item Αναγνωριστικό \textlatin{UUID} έναντι \textlatin{EORI}:
    Ο \textlatin{Keyrock} αποθηκεύει και αναζητά πολιτικές βάσει εσωτερικού
    \textlatin{UUID} και όχι βάσει \textlatin{EORI}, γεγονός που απαιτεί
    προσοχή κατά τη δημιουργία πολιτικών.
    \item Επέκταση \textlatin{keyUsage} πιστοποιητικών: Η επέκταση
    \textlatin{digitalSignature} στα πιστοποιητικά \textlatin{X.509} είναι
    απαραίτητη για τη σωστή λειτουργία του \textlatin{Keyrock}, αλλά
    συχνά παραλείπεται από τα προεπιλεγμένα εργαλεία δημιουργίας
    πιστοποιητικών.
    \item Αρχική καταχώρηση καταναλωτή: Ο \textlatin{Keyrock} δεν
    δημιουργεί αυτόματα χρήστη για εξωτερικούς συμμετέχοντες
    \textlatin{iSHARE} μέχρι την πρώτη τους επικοινωνία, γεγονός που
    απαιτεί συγκεκριμένη σειρά εκκίνησης.
\end{itemize}

% ------------------------------------------------------------------
\section{Μελλοντικές Επεκτάσεις}
\label{sec:future-work}
% ------------------------------------------------------------------

Βάσει των εμπειριών από την υλοποίηση, προτείνονται οι ακόλουθες
κατευθύνσεις για μελλοντική επέκταση:

Η υλοποίηση της ροής \textlatin{H2M} θα επέτρεπε σε χρήστες να αποκτούν
πρόσβαση στα δεδομένα μέσω \textlatin{browser}, χρησιμοποιώντας την
\textlatin{Authorization Code} ροή του \textlatin{OAuth 2.0} και τη διεπαφή
\textlatin{web} του \textlatin{Keyrock} για ταυτοποίηση.

Η μεταφορά του συστήματος σε περιβάλλον \textlatin{Kubernetes} με χρήση των
\textlatin{Helm charts} που παρέχει η \textlatin{FIWARE Foundation} θα
επέτρεπε κλιμάκωση, υψηλή διαθεσιμότητα και δυνατότητα δημόσιας ανάπτυξης.
Αυτό θα καθιστούσε το \textlatin{Data Space} πραγματικά ομοσπονδιακό, με
δυνατότητα σύνδεσης με άλλους οργανισμούς μέσω του πραγματικού
\textlatin{iSHARE Satellite}.

Η νεότερη αρχιτεκτονική του \textlatin{FIWARE Data Space Connector}
βασίζεται σε \textlatin{Verifiable Credentials} (\textlatin{OID4VC}) αντί
για \textlatin{iSHARE JWT}. Η μετάβαση σε αυτό το μοντέλο, ή η υβριδική
υλοποίηση και των δύο, θα αποτελούσε ενδιαφέρον αντικείμενο μελέτης.

Η επέκταση του συστήματος ώστε να υποστηρίζει πολλαπλούς ανεξάρτητους
παρόχους και καταναλωτές θα αποτελούσε ένα πιο ρεαλιστικό σενάριο
\textlatin{Data Space}. Η αρχιτεκτονική που υλοποιήθηκε είναι σχεδιασμένη
να κλιμακωθεί σε αυτή την κατεύθυνση, καθώς κάθε οργανισμός μπορεί να
εκτελεί τον δικό του \textlatin{Kong} και \textlatin{Orion-LD}, με κοινό
μόνο τον \textlatin{iSHARE Satellite} ως αρχή εμπιστοσύνης.

Η προσθήκη μηχανισμού ανακάλυψης δεδομένων (\textlatin{data discovery}),
όπου οι καταναλωτές μπορούν να αναζητούν τα διαθέσιμα \textlatin{datasets}
και να ζητούν πρόσβαση δυναμικά, θα ολοκλήρωνε τον κύκλο ζωής ενός
πλήρους \textlatin{Data Space}.

